\documentclass[10pt,a4paper]{article}
\usepackage{amsmath,amssymb,amsthm,graphicx}
% Card-local minimal preamble (journal unit sigma-2023-014-a-generalization-of-zwegers-mu-function-acco, P2-ahl16 precedent).
%% Verbatim from the pinned source: the author theorem environments and macro block,
%% so every macro and environment of the counted spans is available.
% Omitted and why: the documentclass line (the journal class sigma is not on the host;
%% the assembler supplies the class) and page-geometry/font blocks (they fight the
%% forced xelatex engine).
\usepackage{amsmath,amssymb,amsthm,mathtools}

\numberwithin{equation}{section}
\newtheorem{thm}{Theorem}[section]
\newtheorem{prop}[thm]{Proposition}
\newtheorem{lem}[thm]{Lemma}
\newtheorem{cor}[thm]{Corollary}
\newtheorem{fact}[thm]{Fact}
\newtheorem{conj}[thm]{Conjecture}
\newtheorem{prob}[thm]{Problem}
\newtheorem{ques}[thm]{Question}
\theoremstyle{definition}
\newtheorem{dfn}[thm]{Definition}
\newtheorem{rmk}[thm]{Remark}
\newtheorem{exa}[thm]{Example}
\newcommand{\Hz}{\mathbb{H}}
\newcommand{\Zz}{\mathbb{Z}}
\newcommand{\Cz}{\mathbb{C}}
\newcommand{\hyper}[4]{\left(\begin{matrix} #1 \\ #2 \end{matrix}; #3, #4\right)}
\newcommand{\sn}{\mathrm{sn}\,}
\newcommand{\cn}{\mathrm{cn}\,}
\newcommand{\dn}{\mathrm{dn}\,}
\usepackage{cleveref}
\crefname{originalequation}{original equation}{original equations}
\Crefname{originalequation}{Original equation}{Original equations}
\makeatletter
\newlabel{amuconnect}{{2.2}{source}{}{sourcepointer-1}{}}
\newlabel{amuconnect@cref}{{[originalequation][0][]2.2}{[1][0][]source}{}{}{}}
\makeatother
\title{The f0(x, lambda-prime) connection formula for the generalized Zwegers mu-function}
\author{Shibukawa, Genki, Tsuchimi, Satoshi}
\date{doi:10.3842/SIGMA.2023.014; CC BY 4.0}
\begin{document}
\maketitle
\setcounter{section}{2}
\setcounter{thm}{2}
\setcounter{equation}{3}
\noindent\textit{Editorial note.} All delivered mathematical text is the author's own, copied from the pinned journal source file sources/sigma-2023-014-a-generalization-of-zwegers-mu-function-acco/pinned-original-source.tex (A Generalization of Zwegers' mu-Function According to the q-Hermite-Weber Difference Equation); no statement or proof step was rewritten, reordered or removed.\par\medskip
\noindent\textit{Editorial note.} The counted proof is the author's complete argument: the author substitutes lambda -> lambda-prime in the connection formula amuconnect (2.2), eliminates g0(x) from the resulting two-by-two system eq:different variable (2.5), and closes the remaining theta-ratio identity by the pole-residue constancy argument for the auxiliary function C(x) (residue at x=1 vanishes, C(xq)=C(x), evaluate at x=-a/lambda). In-span labels render from the card's own counters: (2.4) eq:tuchimi connection (statement equation, label line 790; the first row of its gather is nonumber and renders unnumbered exactly as in the source) and (2.5) eq:different variable (proof equation, line 799). Out-of-span navigation pointer declared as external reference: amuconnect (2.2, line 675, used at line 797). The proof tail carries the author's own \textbackslash{}tag*\{\textbackslash{}qed\} plus \textbackslash{}renewcommand\{\textbackslash{}qed\}\{\} device (source lines 820-821), copied verbatim. Seed convention: counters stepped from the source itself and cross-validated against a full host compile of the pinned source (aux ground truth): section 2 of the shared [thm] counter holds lem:lemma1 (2.1) and lem:lemma2 (2.2) before this theorem, and section 2 holds numbered equations (2.1) eq:(26), (2.2) amuconnect and (2.3) eq:prot connec before the span, so the card prints Theorem 2.3 with in-span equations (2.4) and (2.5).\par\medskip
\noindent\textit{Editorial note.} The label thm:tuchimi connection names the counted statement, and the proof environment immediately following it in the source is delivered verbatim as its proof.\par\medskip
\noindent\textit{Editorial note.} The counted statement is carried by the author's own thm environment in the article source.\par\medskip
\noindent\textit{Editorial note.} Journal-unit assembly: card-local preamble (author theorem environments and macros verbatim), article class, no bibliography in the counted span (the span carries no \textbackslash{}cite commands).\par\medskip
\noindent\textit{Editorial scope.} the counted statement and its own proof environment only\par\medskip
\noindent\textit{Allowed standard prerequisites (editorial).}\begin{itemize}
\item The connection formula amuconnect (2.2), source line 675, for f0(x, lambda) in terms of f-infinity(x) and g-infinity(x) with theta-q-function coefficients; the theorem's proof substitutes lambda -> lambda-prime into it.
\item The Jacobi theta function theta\_q and the q-Pochhammer symbols (a)-infinity, (q)-infinity appearing in the theta-coefficients; the proof uses only the pole positions of the auxiliary function C(x) and its periodicity C(xq)=C(x), both stated inside the span (source line 812).
\item The elementary elliptic-function argument that a q-periodic meromorphic function whose prescribed-pole residues vanish is constant, with the constant then read off at one further point (source lines 812-821).
\end{itemize}\medskip
\ifdefined\hypertarget\hypertarget{sourcepointer-1}{}\fi
\noindent\textit{Source reference.} Original equation 2.2: sources/sigma-2023-014-a-generalization-of-zwegers-mu-function-acco/pinned-original-source.tex, source lines 675--675. Navigation pointer invoked at line 797: the counted text cites that numbered equation of the same article, which is not retained in this card, so the reference is declared as a pointer carrying the number the article itself prints for it; the cited equation is an allowed input and none of it is reproduced inside the span.\par\medskip
\begin{thm}
The following equation holds:
\label{thm:tuchimi connection}
\begin{gather}
f_{0}(x,\lambda^{\prime})
 =
 \frac{\theta_q(\lambda^{\prime})\theta_q(ax/\lambda^{\prime})x^\alpha}{\theta_q(\lambda^{\prime}/a)\theta_q(x/\lambda^{\prime})a}f_{\infty}(x,\lambda) \nonumber \\
\label{eq:tuchimi connection}
\hphantom{f_{0}(x,\lambda^{\prime})=}{}
 -
 (a)_\infty(q)_\infty^2\frac{\theta_q(-\lambda^{\prime}/x\lambda)\theta_q(-x)\theta_q(-\lambda\lambda^{\prime}/a)}{\theta_q(x\lambda)\theta_q(\lambda^{\prime}/x)\theta_q(\lambda^{\prime}/a)\theta_q(a/\lambda)}g_{\infty }(x).
\end{gather}
\end{thm}

\begin{proof}
We replace the variable $\lambda$ by $\lambda'$ in the connection formula~(\ref{amuconnect}) of the case of $f_0(x,\lambda)$:
\begin{gather}
\label{eq:different variable}
\begin{pmatrix}f_0(x,\lambda')\\ f_{\infty }(x,\lambda) \end{pmatrix}=-\frac{(q)_\infty}{(q/a)_\infty}\begin{pmatrix}\dfrac{\theta_q(\lambda')\theta_q(ax/\lambda')x^\alpha}{\theta_q(\lambda'/a)\theta_q(x/\lambda')a}&1\\ 1&\dfrac{\theta_q(\lambda)\theta_q(x\lambda/a)}{\theta_q(\lambda/a)\theta_q(x\lambda)}x^{-\alpha}\end{pmatrix}\begin{pmatrix}g_0(x)\\g_\infty(x)\end{pmatrix}.
\end{gather}
By erasing the function $g_0(x)$ in the above (\ref{eq:different variable}), we have
\begin{gather*}
f_0(x,\lambda^{\prime})-\frac{\theta_q(\lambda^{\prime})\theta_q(ax/\lambda^{\prime})x^\alpha}{\theta_q(\lambda^{\prime}/a)\theta_q(x/\lambda^{\prime})a}f_{\infty}(x,\lambda)\\
\qquad=\frac{(q)_\infty}{(q/a)_\infty}\left\{\frac{\theta_q(\lambda)\theta_q(\lambda^{\prime})\theta_q(x\lambda/a)\theta_q(ax/\lambda^{\prime})}{a\theta_q(\lambda/a)\theta_q(\lambda^{\prime}/a)\theta_q(x\lambda)\theta_q(x/\lambda^{\prime})}-1\right\}g_\infty(x)\\
\qquad =\frac{(a)_\infty(q)_\infty^2}{\theta_q(-a)}\frac{\theta_q(-\lambda^{\prime}/x\lambda)\theta_q(-x)}{\theta_q(x\lambda)\theta_q(\lambda^{\prime}/x)}C(x)g_\infty(x),
\end{gather*}
where
\[
C(x)=\frac{\lambda^{\prime}}{ax}\frac{\theta_q(\lambda)\theta_q(\lambda^{\prime})\theta_q(x\lambda/a)\theta_q(ax/\lambda^{\prime})} {\theta_q(-\lambda^{\prime}/x\lambda)\theta_q(-x)\theta_q(\lambda/a)\theta_q(\lambda^{\prime}/a)}-\frac{\theta_q(x\lambda) \theta_q(\lambda^{\prime}/x)}{\theta_q(-\lambda^{\prime}/x\lambda)\theta_q(-x)}.
\]
The function $C(x)$ has a simple pole in the points $x=q^j,\frac{\lambda^{\prime}}{\lambda}q^j$ for some $j\in\Zz$, and satisfies $C(xq)=C(x)$. We calculate the residue at the pole $x=1$,
\begin{align*}
\lim_{x\to1}\theta_q(-x)C(x)&=\frac{\lambda^{\prime}}{a}\frac{\theta_q(\lambda)\theta_q (\lambda^{\prime})\theta_q(\lambda/a)\theta_q(a/\lambda^{\prime})}{\theta_q(-\lambda^{\prime}/\lambda)\theta_q(\lambda/a)\theta_q(\lambda^{\prime}/a)} -\frac{\theta_q(\lambda)\theta_q(\lambda^{\prime})}{\theta_q(-\lambda^{\prime}/\lambda)}\\
&=\frac{\theta_q(\lambda)\theta_q(\lambda^{\prime})}{\theta_q(-\lambda^{\prime}/\lambda)}\left\{\frac{\lambda^{\prime}}{a} \frac{\theta_q(a/\lambda^{\prime})}{\theta_q(\lambda^{\prime}/a)}-1\right\}=0.
\end{align*}
It follows that the doubly periodic function $C(x)$ is constant in $x$.
Hence, we have
\begin{gather*}
C(x)=C(-a/\lambda)=-\frac{\theta_q(-a)\theta_q(-\lambda\lambda^{\prime}/a)}{\theta_q(\lambda^{\prime}/a)\theta_q(a/\lambda)}.\tag*{\qed}
\end{gather*}\renewcommand{\qed}{}
\end{proof}

\end{document}
